\section{Related Work}
\label{sec:related}

\paragraph{Supervised LLM routing.}
RouteLLM learns routers from human preference data and augmented pairwise labels~\citep{ong2024routellm}.
HybridLLM trains a query encoder with supervision from model quality scores~\citep{ding2024hybrid}.
GraphRouter routes among models and agents with graph-based supervised training~\citep{feng2025graphrouter}.
AgentRouter formulates multi-agent QA as knowledge-graph--guided routing, training a heterogeneous GNN with soft supervision from per-agent performance~\citep{zhang2026agentrouter}.
RouterBench provides benchmarks for comparing supervised routers~\citep{hu2024routerbench}.
These systems \emph{learn} routing inputs from preference data, outcome labels, or structured supervision; our contrast is at \textbf{signal extraction}: we pre-specify unsupervised statistics and ask whether they predict routing need before any routing-specific representation learning.

\paragraph{Cost-aware and cascade routing.}
FrugalGPT and related work cascade or select models using feedback from generations or external scorers~\citep{chen2024frugalgpt}.
Smoothie combines multiple model outputs with label-free weighting over candidates~\citep{guha2024smoothie}.
Such methods often require post-generation information, multiple samples, or reward models; we focus on \textbf{pre-decision signals} from query text and model responses in a fixed capability-differentiated pool, and we report accuracy--cost \textbf{Pareto} curves rather than accuracy alone.

\paragraph{Uncertainty and confidence signals.}
Work on calibration and self-evaluation shows that LLM probabilities can relate to correctness~\citep{kadavath2022know,plaut2024calibration,he2023calibration}.
Semantic uncertainty uses multiple samples to estimate confidence~\citep{kuhn2023semantic}.
These lines motivate \textbf{model-response} and \textbf{cross-model comparative} signals but do not systematically test nested unsupervised signal types against \textbf{routing need} $r(q)$---$\Mlo$ fails, $\Mhi$ succeeds---nor separate signal informativeness (H1--H3) from deployable routing quality (H4).

\paragraph{Positioning.}
Table~\ref{tab:contrast} summarizes how our setting differs from typical supervised routers.
Prior routing is largely \textbf{supervised at the router}: labels or preferences \emph{define} what the router measures and predicts.
We ask whether \textbf{unsupervised signals}---query-derived, model-response, and cross-model comparative---carry information about $r(q)$ and whether simple threshold policies built from them improve the accuracy--cost trade-off relative to always using one model.
We cite supervised systems as field context; reproducing them is out of scope (Section~\ref{sec:discussion}).

\begin{table}[t]
  \centering
  \small
  \begin{tabular}{@{}p{0.42\linewidth}p{0.48\linewidth}@{}}
    \toprule
    \textbf{Typical supervised router} & \textbf{This work} \\
    \midrule
    Router input: classifier on query ($+$ history), trained on prefs or outcomes &
    \textbf{Unsupervised signals} $\rightarrow$ feature vector $x(q)$ \\
    Labels \textbf{define} the router end-to-end &
    Labels used \textbf{only on calib}: signal analysis + $(\lambda,\tau)$ fit \\
    Signal source: learned representations from routing data &
    Query-derived, model-response, cross-model comparative (Section~\ref{sec:signal-extraction}) \\
    Evaluation: routing accuracy or win rate &
    \textbf{Pareto} over task $\mathrm{Acc}(\pi)$ and $\mathrm{Cost}(\pi)$ on test \\
    \bottomrule
  \end{tabular}
  \caption{Contrast with supervised LLM routing (program §2). Unsupervised describes layer~1 signal extraction; calib uses offline $r(q)$ for analysis and policy fit only.}
  \label{tab:contrast}
\end{table}
